\documentclass[12pt,a4paper]{article}

% ── Encodage & langue ─────────────────────────────────────
\usepackage[utf8]{inputenc}
\usepackage[T1]{fontenc}
\usepackage[french]{babel}

% ── Mise en page ──────────────────────────────────────────
\usepackage[top=2.5cm,bottom=2.5cm,left=3cm,right=3cm]{geometry}
\usepackage{setspace}
\onehalfspacing

% ── Mathématiques ─────────────────────────────────────────
\usepackage{amsmath,amssymb,amsthm}
\usepackage{mathtools}

% ── Graphiques & couleurs ─────────────────────────────────
\usepackage{graphicx}
\usepackage[dvipsnames,svgnames,x11names]{xcolor}
\usepackage{tikz}
\usetikzlibrary{arrows.meta,positioning,shapes,fit,backgrounds}

% ── Tableaux ──────────────────────────────────────────────
\usepackage{booktabs}
\usepackage{array}
\usepackage{multirow}

% ── Liens & métadonnées ───────────────────────────────────
\usepackage{hyperref}
\hypersetup{
  colorlinks=true,
  linkcolor=NavyBlue,
  citecolor=NavyBlue,
  urlcolor=NavyBlue,
  pdftitle={CommunityTrust Protocol},
  pdfauthor={OpenScience Community},
}

% ── Algorithmes ───────────────────────────────────────────
\usepackage{algorithm}
\usepackage{algpseudocode}

% ── En-têtes ──────────────────────────────────────────────
\usepackage{fancyhdr}
\pagestyle{fancy}
\fancyhf{}
\rhead{\small\textcolor{gray}{CTP v0.1 — OpenScience Community}}
\lhead{\small\textcolor{gray}{CommunityTrust Protocol}}
\cfoot{\thepage}

% ── Environnements théoriques ─────────────────────────────
\theoremstyle{definition}
\newtheorem{definition}{Définition}[section]
\newtheorem{exemple}{Exemple}[section]

\theoremstyle{plain}
\newtheorem{theoreme}{Théorème}[section]
\newtheorem{proposition}{Proposition}[section]
\newtheorem{lemme}{Lemme}[section]
\newtheorem{corollaire}{Corollaire}[section]

\theoremstyle{remark}
\newtheorem{remarque}{Remarque}[section]

% ── Commandes personnalisées ──────────────────────────────
\newcommand{\CTP}{\textsc{ctp}}
\newcommand{\trustscore}{\tau}
\newcommand{\R}{\mathbb{R}}
\newcommand{\N}{\mathbb{N}}
\newcommand{\E}{\mathbb{E}}
\newcommand{\1}{\mathbf{1}}
\newcommand{\clip}{\operatorname{clip}}

% ─────────────────────────────────────────────────────────
\begin{document}

% ── Page de titre ─────────────────────────────────────────
\begin{titlepage}
\centering
\vspace*{1.5cm}

{\Large \textsc{OpenScience Community}}\\[0.4cm]
{\small Protocole Technique — Document de Référence}\\[2cm]

\hrule height 0.5pt
\vspace{0.6cm}

{\Huge \textbf{CommunityTrust Protocol}}\\[0.5cm]
{\LARGE \textit{CTP v1.1}}\\[0.3cm]

\vspace{0.4cm}
\hrule height 0.5pt
\vspace{1.2cm}

{\large \textit{``Je suis parce que nous sommes''}}\\[0.2cm]
{\small \textit{--- Philosophie Ubuntu, Afrique australe}}

\vspace{1.5cm}

{\large Un protocole universel de confiance communautaire\\
inspiré des tontines africaines,\\
formalisé par les mathématiques,\\
validé par simulation,\\
conçu pour les smart contracts.}

\vspace{1.5cm}

\begin{tabular}{rl}
\textbf{Auteur}       & Ange AWALA \\
\textbf{Affiliation}  & OpenScience Community \\
\textbf{Contact}      & openscience.community@proton.me \\[0.3cm]
\textbf{Version}      & 1.1 --- Corrections et extensions \\
\textbf{Statut}       & Recherche ouverte \\
\textbf{Licence}      & Creative Commons BY-SA 4.0 \\
\textbf{Dépôt}        & Zenodo --- DOI à venir \\
\end{tabular}

\vfill
{\small \today}
\end{titlepage}

% ── Résumé ────────────────────────────────────────────────
\begin{abstract}
\noindent
Les systèmes de confiance existants mesurent la popularité, la richesse ou
le statut institutionnel. Le \textbf{CommunityTrust Protocol} (\CTP{}) propose
une alternative radicale~: mesurer la \emph{fiabilité comportementale dans le
temps}, indépendamment de tout capital financier ou social préexistant.

Inspiré de la logique des tontines africaines --- systèmes de financement
communautaire décentralisé où la confiance collective remplace la garantie
bancaire --- \CTP{} formalise un moteur de réputation universel applicable à
tout système communautaire~: réseaux sociaux, organisations autonomes
décentralisées (DAOs), coopératives de connaissance, systèmes de
micro-finance.

Ce document présente~: (1) la spécification conceptuelle complète du
protocole, (2) le modèle mathématique rigoureux incluant la dynamique du
\emph{TrustScore}, le mécanisme de responsabilité des \emph{Witnesses} et la
hiérarchie des \emph{Nodes}, (3) une analyse par la théorie des jeux
garantissant la résistance aux comportements stratégiques, et (4) une
validation empirique par simulation Monte Carlo sur 52 périodes.

Les résultats démontrent que les comportements honnêtes dominent strictement
les stratégies de free-riding, collusion et désertion sous les paramètres
recommandés. Les trois propositions théoriques sont vérifiées empiriquement.

\medskip\noindent
\textbf{Nouveautés v1.1~:} (a) Capital de Pardon --- taux de gain adaptatif
$\eta_i(t)$ selon l'ancienneté~; (b) similarité sémantique directionnelle
$\operatorname{sim}(j \to j')$ calculable on-chain~; (c) analyse de
complexité algorithmique et coûts gas~; (d) référence fondatrice de Marsh
(1994) positionnée~; (e) section Limites explicites~; (f) correction du
paramètre $\rho_r$ dans le module Finance.

\medskip
\noindent\textbf{Mots-clés~:} confiance communautaire, tontine, réputation
décentralisée, théorie des jeux, smart contracts, science ouverte, protocole
P2P.
\end{abstract}

\newpage
\tableofcontents
\newpage

% ═════════════════════════════════════════════════════════
\section{Introduction}
% ═════════════════════════════════════════════════════════

\subsection{Motivation}

Les communautés humaines ont développé depuis des millénaires des mécanismes
sophistiqués de confiance collective sans recours aux institutions financières
formelles. La \textbf{tontine} --- pratique répandue à travers l'Afrique,
l'Asie et la diaspora africaine mondiale --- en constitue l'exemple le plus
remarquable~: un groupe de personnes contribue périodiquement à un pot commun,
redistribué à tour de rôle selon des règles négociées collectivement. Ce
système fonctionne sans contrat légal, sans garantie bancaire, uniquement sur
la base de la réputation et de la pression sociale.

Ce qui rend les tontines efficaces n'est pas l'argent en jeu, mais
l'\textbf{information communautaire}~: chaque membre connaît le comportement
passé de ses pairs, et cette connaissance structure les incitations présentes
et futures. La confiance est une ressource circulante, construite lentement et
perdue rapidement.

Les systèmes numériques contemporains ont largement échoué à capturer cette
logique. Les réseaux sociaux mesurent la popularité algorithmique. Les DAOs
confondent richesse et gouvernance. Les plateformes de micro-finance importent
des modèles bancaires inadaptés aux contextes de faible bancarisation. Aucun
ne formalise ce que les tontines font naturellement~: \emph{quantifier la
fiabilité relationnelle dans le temps}.

\subsection{Contribution}

Le \textbf{CommunityTrust Protocol} (\CTP{}) apporte les contributions
suivantes~:

\begin{enumerate}
  \item \textbf{Universalité}~: un moteur de réputation unique applicable à
    tout système communautaire, paramétrable sans modification du code central.

  \item \textbf{Ancrage théorique}~: une formalisation mathématique rigoureuse
    des mécanismes de confiance communautaire africaine, jusqu'ici transmis
    oralement.

  \item \textbf{Résistance stratégique}~: une analyse par la théorie des jeux
    prouvant que les comportements honnêtes constituent un équilibre dominant
    sous les paramètres recommandés.

  \item \textbf{Implémentabilité}~: une architecture directement traduisible
    en smart contracts, conçue dès l'origine pour l'exécution décentralisée.

  \item \textbf{Hiérarchie de Nodes}~: un système multi-niveaux
    (Meta-Node~/ Node~/ Sub-Node) permettant la coexistence de communautés
    autonomes partageant un socle commun de règles éthiques.
\end{enumerate}

\subsection{Structure du document}

La Section~\ref{sec:concepts} introduit le vocabulaire fondateur. La
Section~\ref{sec:modele} présente le modèle mathématique complet. La
Section~\ref{sec:jeux} analyse les équilibres stratégiques. La
Section~\ref{sec:simulation} valide les résultats par simulation. La
Section~\ref{sec:architecture} décrit l'architecture des Nodes. La
Section~\ref{sec:ethique} pose les principes éthiques fondateurs. La
Section~\ref{sec:conclusion} conclut et ouvre les perspectives.

% ═════════════════════════════════════════════════════════
\section{Spécification Conceptuelle}
\label{sec:concepts}
% ═════════════════════════════════════════════════════════

\subsection{Vocabulaire fondateur}

\CTP{} repose sur cinq concepts primitifs dont les définitions sont
immuables~: aucune implémentation ne peut les redéfinir.

\begin{definition}[Signal]
Un \textbf{Signal} $s$ est tout acte brut produit librement par un membre~:
publication, paiement, vote, partage de ressource. Le Signal n'a pas de valeur
a priori --- il existe simplement comme événement observable.
\end{definition}

\begin{definition}[Contribution]
Une \textbf{Contribution} est un Signal validé collectivement par un nombre
suffisant de Witnesses. C'est l'unité de base du protocole. Sans validation
communautaire, il n'existe pas de Contribution.
\end{definition}

\begin{definition}[Witness]
Un \textbf{Witness} est un membre dont le TrustScore dépasse un seuil fixé par
le Node, habilité à valider des Signals. Le Witness engage sa propre réputation
dans l'acte de validation~: valider un Signal de mauvaise qualité a un coût.
\end{definition}

\begin{definition}[TrustScore]
Le \textbf{TrustScore} $\trustscore_{i,j}(t) \in [0,1]$ est le score de
réputation du membre $i$ dans le Node $j$ au temps $t$. Il ne se vend pas, ne
se transfère pas et ne s'achète pas. Il se construit uniquement par le
comportement observé dans le temps.
\end{definition}

\begin{definition}[Node]
Un \textbf{Node} $N_j$ est une communauté qui implémente \CTP{}. Il définit
ses propres paramètres de validation dans les limites fixées par le Meta-Node.
Des exemples de Nodes~: une tontine de quartier, un groupe de recherche
scientifique, une coopérative agricole, une DAO.
\end{definition}

\subsection{Cycle de vie d'une Contribution}

Le cycle de vie complet d'un Signal est représenté à la
Figure~\ref{fig:cycle}.

\begin{figure}[h]
\centering
\begin{tikzpicture}[
  node distance=1.3cm and 2.5cm,
  box/.style={rectangle,rounded corners=4pt,draw=NavyBlue!70,
              fill=NavyBlue!8,minimum width=3.2cm,minimum height=0.9cm,
              font=\small,align=center},
  decision/.style={diamond,draw=Orange!80,fill=Orange!10,
                   minimum width=2.4cm,minimum height=0.9cm,
                   font=\small,align=center,aspect=2},
  arr/.style={-{Stealth[length=5pt]},NavyBlue!70,thick},
  warn/.style={-{Stealth[length=5pt]},Red!70,thick,dashed},
]

\node[box] (prod) {Membre produit\\un Signal $s$};
\node[box,right=of prod] (soum) {Soumission\\au Node $N_j$};
\node[decision,right=of soum] (wit) {$k$ Witnesses\\disponibles~?};
\node[box,below=of wit] (rej1) {Signal rejeté\\(communauté trop petite)};
\node[decision,right=of wit] (val) {Seuil de\\validation atteint~?};
\node[box,below=of val] (rej2) {Signal non validé\\(qualité insuffisante)};
\node[box,right=of val] (contrib) {Signal devient\\Contribution};
\node[box,right=of contrib] (update) {Mise à jour\\TrustScores};
\node[decision,below=of update] (audit) {Audit~?};
\node[box,left=of audit] (cout) {Coût Witnesses\\$\phi\cdot\trustscore_w$};

\draw[arr] (prod)    -- (soum);
\draw[arr] (soum)    -- (wit);
\draw[arr] (wit)     -- node[right,font=\tiny]{non} (rej1);
\draw[arr] (wit)     -- node[above,font=\tiny]{oui} (val);
\draw[arr] (val)     -- node[right,font=\tiny]{non} (rej2);
\draw[arr] (val)     -- node[above,font=\tiny]{oui} (contrib);
\draw[arr] (contrib) -- (update);
\draw[arr] (update)  -- (audit);
\draw[warn](audit)   -- node[above,font=\tiny]{oui} (cout);

\end{tikzpicture}
\caption{Cycle de vie d'un Signal dans \CTP{}.}
\label{fig:cycle}
\end{figure}

\subsection{Principe fondateur}

Le principe central de \CTP{} peut s'énoncer ainsi~:

\begin{center}
\fbox{\parbox{0.85\textwidth}{\centering\itshape
La confiance d'un membre est la somme de ses comportements observés,\\
pondérée par la fiabilité de ceux qui les ont observés.
}}
\end{center}

Ce principe implique une récursivité constitutive~: la valeur d'une validation
dépend de la réputation du validateur, qui elle-même dépend de ses validations
passées. Cela crée une structure d'incitations auto-renforçante favorable aux
comportements honnêtes.

% ═════════════════════════════════════════════════════════
\section{Modèle Mathématique}
\label{sec:modele}
% ═════════════════════════════════════════════════════════

\subsection{Espaces et notations}

Soit un Meta-Node $\mathcal{M}$ contenant $m$ Nodes $\mathcal{N} = \{N_1,
\ldots, N_m\}$. Chaque Node $N_j$ contient un ensemble de membres
$\mathcal{M}_j = \{1, \ldots, n_j\}$. Le système évolue en périodes discrètes
$t \in \N$, dont la durée (hebdomadaire, mensuelle) est définie par chaque
Node.

\subsection{Dynamique du TrustScore}

\subsubsection{Équation d'évolution}

Le TrustScore du membre $i$ dans le Node $j$ évolue selon~:

\begin{equation}
\boxed{
\trustscore_{i,j}(t+1) = \clip\!\Bigl(
  \trustscore_{i,j}(t) + G_{i,j}(t) - L_{i,j}(t) - D_{i,j}(t) - W_{i,j}(t),\;
  0,\; 1
\Bigr)
}
\label{eq:trust}
\end{equation}

où $\clip(x, 0, 1) = \max(0, \min(1, x))$ garantit la bornitude dans $[0,1]$.

\subsubsection{Force 1 --- Gain par contribution}

\begin{equation}
G_{i,j}(t) = \eta \sum_{k=1}^{K_{i,j}(t)} q(s_k)
\label{eq:gain}
\end{equation}

où $\eta \in (0,1)$ est le taux d'apprentissage du Node (petit, car la
confiance se construit lentement), $K_{i,j}(t)$ est le nombre de Contributions
validées à la période $t$, et $q(s_k) \in [0,1]$ est la qualité de la
$k$-ième Contribution définie à l'équation~\eqref{eq:quality}.

\subsubsection{Force 2 --- Perte par défaut}

\begin{equation}
L_{i,j}(t) = \kappa \cdot \1_{\{\text{défaut}_i\}}(t)
\label{eq:loss}
\end{equation}

où $\kappa \in (0,1)$ est la pénalité de défaut, avec la contrainte
$\kappa \gg \eta$~: on perd rapidement ce que l'on a construit lentement,
conformément à la dynamique réelle de la confiance humaine.

\subsubsection{Force 3 --- Dépréciation par inactivité}

\begin{equation}
D_{i,j}(t) = \delta \cdot \trustscore_{i,j}(t) \cdot \1_{\{\text{inactif}_i\}}(t)
\label{eq:depreciation}
\end{equation}

où $\delta \in (0,1)$ est le taux de dépréciation. La dépréciation est
proportionnelle au score courant~: les membres établis ont davantage à perdre
en restant inactifs. Ce mécanisme empêche de vivre sur ses acquis.

\subsubsection{Extension v1.1 --- Capital de Pardon}

Un défaut accidentel d'un membre de longue date ne doit pas être traité
de façon identique à une fraude d'un nouveau venu. \CTP{} v1.1 introduit
un \textbf{taux de gain adaptatif}~:

\begin{equation}
\boxed{\eta_i(t) = \eta \cdot \left(1 + \alpha_P \cdot
\frac{h_{i,j}(t)}{T_{\text{ref}}}\right)}
\label{eq:pardon}
\end{equation}

où $\alpha_P \in [0,1]$ est l'amplificateur de pardon et
$T_{\text{ref}}$ est une période de référence (valeur recommandée~:
52 périodes). Un membre actif depuis deux ans ($h = 104$) remonte
son score environ deux fois plus vite qu'un nouveau venu après un défaut.

\begin{remarque}
Le Capital de Pardon ne réduit pas la pénalité $\kappa$ --- la chute
reste aussi brutale pour tout le monde. Seule la \emph{vitesse de
reconstruction} est modulée par l'ancienneté. L'équité de la punition
est préservée~; la reconnaissance de l'historique ne l'est pas.
\end{remarque}

\begin{proposition}[Convergence avec Capital de Pardon]
Sous comportement constant et Capital de Pardon actif, le TrustScore
d'un membre de longue date converge vers~:
$$\trustscore_i^* = \frac{\eta_i \cdot \bar{q}}{\eta_i \cdot \bar{q} + \epsilon}$$
qui est strictement supérieur à $\trustscore^*$ sans Capital de Pardon,
toutes choses égales par ailleurs.
\end{proposition}

\subsubsection{Force 4 --- Coût de mauvaise validation}

\begin{equation}
W_{i,j}(t) = \phi \sum_{s \,:\, i \in \mathcal{W}(s),\; s \text{ audité}}
\frac{\trustscore_{i,j}(t)}{|\mathcal{W}(s)|}
\label{eq:witness_cost}
\end{equation}

où $\phi \in (0,1)$ est le coefficient de responsabilité et $\mathcal{W}(s)$
l'ensemble des Witnesses du Signal $s$. Ce terme est la contribution centrale
de \CTP{} par rapport aux systèmes existants~: valider un Signal de mauvaise
qualité a un coût réel, partagé entre les Witnesses et proportionnel à leur
réputation.

\subsection{Qualité d'une Contribution}

La qualité d'un Signal validé est~:

\begin{equation}
q(s) = \bar{\trustscore}_{\mathcal{W}(s)} \cdot \rho(s)
= \frac{\sum_{w \in \mathcal{W}(s)} \trustscore_{w,j}(t)}{|\mathcal{W}(s)|}
\cdot \min\!\left(\frac{|\mathcal{W}(s)|}{k_j},\, 1\right)
\label{eq:quality}
\end{equation}

où $k_j$ est le seuil minimal de Witnesses du Node $j$, et $\rho(s) \in [0,1]$
est le taux de consensus. Une Contribution validée à l'unanimité par des
membres expérimentés vaut davantage qu'une Contribution validée de justesse par
des débutants.

\subsection{Seuil adaptatif de Witnesses}

Pour éviter le blocage dans les petites communautés, le seuil de validation est
adaptatif~:

\begin{equation}
k_j = \max\!\left(k_{\min},\, \left\lfloor \alpha_k \cdot |\mathcal{M}_j| \right\rfloor\right)
\label{eq:k}
\end{equation}

avec $k_{\min} = 3$ et $\alpha_k = 0{,}20$ (valeurs recommandées par le
Meta-Node). Un Node de 20 membres requiert $\max(3, 4) = 4$ Witnesses~; un
Node de 5 membres requiert $\max(3, 1) = 3$ Witnesses.

\subsection{TrustScore Global et hiérarchie des Nodes}

\subsubsection{Score Global dans le Meta-Node}

Un membre appartenant à $m'$ Nodes simultanément possède un TrustScore Global~:

\begin{equation}
\trustscore_i^{\mathcal{M}}(t)
= \sum_{j=1}^{m'} \omega_{i,j}(t) \cdot \trustscore_{i,j}(t)
\label{eq:global}
\end{equation}

avec les poids~:

\begin{equation}
\omega_{i,j}(t) = \frac{h_{i,j}(t) \cdot \rho_j}{\displaystyle\sum_{j'=1}^{m'} h_{i,j'}(t) \cdot \rho_{j'}}
\label{eq:omega}
\end{equation}

où $h_{i,j}(t)$ est le nombre de périodes actives du membre $i$ dans le Node
$j$, et $\rho_j$ est l'indice de difficulté du Node défini ci-dessous.

\subsubsection{Indice de difficulté d'un Node}

\begin{equation}
\rho_j = \frac{k_j \cdot \bar{\trustscore}_j}{k_{\min} \cdot \trustscore_{\min}}
\label{eq:difficulty}
\end{equation}

où $\bar{\trustscore}_j$ est le TrustScore moyen des membres actifs du Node
$j$, et $\trustscore_{\min}$ est la valeur minimale observée sur l'ensemble du
Meta-Node. \textbf{L'indice de difficulté n'est jamais déclaré par le Node
lui-même~: il est calculé automatiquement à partir de données observables.}

\subsubsection{Portabilité inter-Nodes}

Lorsqu'un membre $i$ rejoint un nouveau Node $j'$ sans historique, son score
initial est~:

\begin{equation}
\trustscore_{i,j'}(0) = \lambda_{\text{eff}}(m_i) \cdot \trustscore_i^{\mathcal{M}}(t)
\label{eq:portability}
\end{equation}

avec le coefficient de portabilité effectif~:

\begin{equation}
\lambda_{\text{eff}}(m_i) = \lambda \cdot \mu^{m_i} \cdot \operatorname{sim}(j_{\text{origine}}, j')
\label{eq:lambda_eff}
\end{equation}

où $m_i$ est le nombre de migrations récentes du membre, $\mu \in (0,1)$ est
la pénalité de migration fréquente, et $\operatorname{sim}(j, j') \in [0,1]$
est la similarité sémantique entre les deux Nodes, définie par le Meta-Node.
Le Node d'accueil peut abaisser $\lambda$ en-deçà de cette valeur, jamais
l'augmenter.

\paragraph{Extension v1.1 --- Similarité directionnelle basée sur les flux.}

En v0.1, $\operatorname{sim}(j, j')$ était spécifiée sans être calculable.
\CTP{} v1.1 propose une formalisation directe, on-chain, basée sur les
comportements réels des membres~:

\begin{equation}
\boxed{\operatorname{sim}(j \to j') =
\frac{|\mathcal{M}_j \cap \mathcal{M}_{j'}|}{|\mathcal{M}_j|}}
\label{eq:sim}
\end{equation}

C'est la proportion des membres du Node $j$ qui appartiennent aussi au
Node $j'$. Cette mesure est~:

\begin{itemize}
  \item \textbf{Directionnelle}~: $\operatorname{sim}(j \to j') \neq
    \operatorname{sim}(j' \to j)$ en général. Si beaucoup de
    mathématiciens rejoignent le Node Physique mais peu de physiciens
    rejoignent Mathématiques, la similarité n'est pas symétrique.
  \item \textbf{Calculable on-chain}~: les ensembles $\mathcal{M}_j$
    et $\mathcal{M}_{j'}$ sont disponibles directement dans CTPCore.
  \item \textbf{Auto-corrective}~: si les habitudes de migration
    changent, $\operatorname{sim}$ change automatiquement sans
    intervention humaine.
  \item \textbf{Résistante à la manipulation}~: un individu seul ne
    peut pas modifier significativement la similarité entre deux
    Nodes de taille raisonnable.
\end{itemize}

\begin{proposition}[Bornitude de sim]
$\operatorname{sim}(j \to j') \in [0,1]$ pour tout $j, j'$, avec
$\operatorname{sim}(j \to j) = 1$ et
$\operatorname{sim}(j \to j') = 0$ si $\mathcal{M}_j \cap \mathcal{M}_{j'} = \emptyset$.
\end{proposition}

\subsection{Pouvoir de vote}

Le droit de décision du membre $i$ dans le Node $j$ est~:

\begin{equation}
V_{i,j}(t) = \frac{\trustscore_{i,j}(t)}{\displaystyle\sum_{k \in \mathcal{M}_j} \trustscore_{k,j}(t)}
\label{eq:vote}
\end{equation}

Le vote est proportionnel au TrustScore, non à la richesse ni à l'ancienneté brute.

\subsection{Propriétés formelles}

\begin{proposition}[Bornitude]
Pour tout membre $i$, Node $j$ et temps $t \geq 0$~:
$$\trustscore_{i,j}(t) \in [0, 1]$$
\end{proposition}

\begin{proof}
Par initialisation $\trustscore_{i,j}(0) \in [0,1]$. La fonction $\clip$ dans
l'équation~\eqref{eq:trust} garantit que la propriété est préservée à chaque
étape par induction.
\end{proof}

\begin{proposition}[Convergence sous comportement constant]
Si un membre contribue régulièrement avec qualité moyenne $\bar{q}$, sans
défaut ($p_i = 0$) et sans inactivité ($a_i = 1$), son TrustScore converge
vers~:
$$\trustscore_i^* = \frac{\eta \cdot \bar{q}}{\eta \cdot \bar{q} + \epsilon}$$
où $\epsilon \to 0$ lorsque $\phi \to 0$. En pratique, $\trustscore_i^*$ tend
vers $1$ pour un membre parfaitement fiable.
\end{proposition}

\begin{proposition}[Résistance à l'accumulation passive]
Un membre inactif voit son TrustScore décroître géométriquement~:
$$\trustscore_{i,j}(t) = \trustscore_{i,j}(0) \cdot (1-\delta)^t \xrightarrow{t \to \infty} 0$$
Aucun membre ne peut maintenir un score élevé sans contribution active.
\end{proposition}

% ═════════════════════════════════════════════════════════
\section{Théorie des Jeux}
\label{sec:jeux}
% ═════════════════════════════════════════════════════════

\subsection{Cadre général}

Chaque membre $i$ est modélisé comme un joueur rationnel maximisant son
utilité espérée sur un horizon $T$~:

\begin{equation}
U_i = \E\left[\sum_{t=0}^{T} \beta^t \cdot u_i(t)\right]
\label{eq:utility}
\end{equation}

où $\beta \in (0,1)$ est le facteur d'escompte temporel --- un membre
\og{}myope\fg{} a $\beta \approx 0$, un membre avec vision à long terme a
$\beta \approx 1$ --- et l'utilité instantanée est~:

\begin{equation}
u_i(t) = f(\trustscore_{i,j}(t)) - c_i(t)
\end{equation}

où $f : [0,1] \to \R_+$ est une fonction croissante représentant les droits
conférés par le TrustScore (accès aux ressources, pouvoir de vote, financement)
et $c_i(t)$ est le coût réel de la contribution à la période $t$.

\subsection{Jeu 1 --- Le Passager Clandestin}

\subsubsection{Formalisation}

Le membre $i$ choisit entre contribuer (stratégie $C$, coût $c > 0$) ou
parasiter (stratégie $P$, coût nul). Sous stratégie $P$, le score décroît
par la dépréciation~: $\trustscore_i(t) = \trustscore_i(0)(1-\delta)^t$.

\subsubsection{Condition de dominance}

La stratégie $C$ domine $P$ si et seulement si~:

\begin{equation}
\sum_{t=0}^{T} \beta^t \left[f(\trustscore_i^C(t)) - f\!\left(\trustscore_i(0)(1-\delta)^t\right)\right] > c \cdot T
\label{eq:freerider}
\end{equation}

\begin{proposition}[Neutralisation du free-riding]
Si $f(\trustscore)$ est suffisamment croissante --- c'est-à-dire si les droits
associés au TrustScore ont une valeur réelle --- et si $\beta$ est
suffisamment proche de 1, alors la stratégie $C$ domine strictement $P$.
\end{proposition}

\begin{remarque}
\textbf{Implication critique pour l'implémentation~:} les droits associés au
TrustScore ne sont pas optionnels. Un Node qui n'attache aucun droit réel au
score crée mécaniquement du free-riding. La valeur de $f$ est le premier levier
de conception.
\end{remarque}

\subsection{Jeu 2 --- La Collusion entre Witnesses}

\subsubsection{Formalisation}

Ce jeu est modélisé comme un \textbf{dilemme du prisonnier répété} entre un
membre $i$ et ses Witnesses complices potentiels. À chaque période, chaque
Witness choisit entre valider honnêtement (H) ou colluder (K).

\subsubsection{Analyse des payoffs}

Pour un Witness $w$ participant à la collusion~:

\begin{align}
G_{\text{collusion}} &= \eta \cdot q_{\text{artificiel}} \cdot N_{\text{validés}} \\
C_{\text{détection}} &= \phi \cdot \trustscore_w(t) \cdot \mathbb{P}(\text{détection})
\end{align}

\begin{proposition}[Condition anti-collusion]
La collusion est irrationnelle pour le Witness si~:

\begin{equation}
\phi > \frac{\eta \cdot q_{\text{artificiel}}}{\trustscore_w(t) \cdot \mathbb{P}(\text{détection})}
\label{eq:anticollusion}
\end{equation}

En particulier, la condition $\phi \geq 3\eta$ est suffisante pour
$\trustscore_w \geq 0{,}3$ et $\mathbb{P}(\text{détection}) \geq 0{,}15$.
\end{proposition}

\subsubsection{Équilibre dans le jeu répété}

Dans un jeu répété à horizon infini avec $\beta$ suffisamment grand, la
stratégie \textit{Tit-for-Tat honnête} --- être honnête, punir immédiatement
toute collusion détectée, reprendre la coopération si le comportement redevient
honnête --- constitue un équilibre de Nash sous~:

\begin{equation}
\beta > \frac{\phi \cdot \trustscore_w \cdot \mathbb{P}(\text{détection})}{\eta \cdot q_{\text{artificiel}}}
\label{eq:tft}
\end{equation}

\subsection{Jeu 3 --- Le Déserteur Stratégique}

\subsubsection{Formalisation}

Le déserteur entre dans un Node, monte son score minimalement, extrait les
droits disponibles, puis migre vers un nouveau Node en emportant sa
portabilité. Son utilité sur $m$ migrations est~:

\begin{equation}
U_{\text{déserteur}} = \sum_{j=1}^{m} \left[f(\lambda_{\text{eff}}(j) \cdot \trustscore^{\mathcal{M}}) - c_{\min} \cdot T_j\right]
\end{equation}

\subsubsection{Mécanismes de neutralisation}

Trois mécanismes combinés rendent cette stratégie irrationnelle~:

\begin{enumerate}
  \item \textbf{Dépréciation du score global~:} $\trustscore^{\mathcal{M}}(t) \to 0$ en l'absence de contribution réelle.
  \item \textbf{Pénalité de migration~:} $\lambda_{\text{eff}}(m) = \lambda \cdot \mu^m \cdot \operatorname{sim}(j, j')$ décroît exponentiellement avec $m$.
  \item \textbf{Distance sémantique~:} $\operatorname{sim}(j, j') \approx 0$ pour des domaines éloignés.
\end{enumerate}

\begin{proposition}[Neutralisation de la désertion]
Pour $\mu \in [0{,}6, 0{,}8]$ et $\beta$ suffisamment grand, la stratégie de
désertion est dominée par la contribution honnête dans un seul Node.
\end{proposition}

\subsection{Synthèse des équilibres}

\begin{table}[h]
\centering
\caption{Résumé des conditions de résistance stratégique de \CTP{}.}
\label{tab:equilibres}
\begin{tabular}{llll}
\toprule
\textbf{Comportement} & \textbf{Mécanisme} & \textbf{Condition} & \textbf{Valeur rec.} \\
\midrule
Passager clandestin & Dépréciation $\delta$ + valeur de $f$ & $f$ croissante réelle & $\delta = 0{,}02$ \\
Collusion Witnesses & Responsabilité $\phi$ & $\phi \geq 3\eta$ & $\phi = 0{,}15$ \\
Déserteur & Pénalité $\mu$ + distance $\operatorname{sim}$ & $\mu \in [0{,}6, 0{,}8]$ & $\mu = 0{,}70$ \\
\bottomrule
\end{tabular}
\end{table}

% ═════════════════════════════════════════════════════════
\section{Validation par Simulation}
\label{sec:simulation}
% ═════════════════════════════════════════════════════════

\subsection{Protocole expérimental}

La simulation est conduite sur un Node de 12 membres répartis en quatre
groupes comportementaux~: 5 membres honnêtes, 3 passagers clandestins, 2
colludeurs, 2 déserteurs. Le Node est simulé sur $T = 52$ périodes
(équivalent d'une année hebdomadaire). Les paramètres utilisés sont ceux du
Tableau~\ref{tab:parametres}.

\begin{table}[h]
\centering
\caption{Paramètres complets de \CTP{} v0.1.}
\label{tab:parametres}
\begin{tabular}{clll}
\toprule
\textbf{Paramètre} & \textbf{Rôle} & \textbf{Valeur} & \textbf{Contrainte} \\
\midrule
$\eta$   & Taux de gain de base           & $0{,}05$  & $\eta \ll \kappa$ \\
$\alpha_P$ & Amplificateur de pardon      & $0{,}50$  & $\alpha_P \in [0,1]$ \\
$T_{\text{ref}}$ & Période de référence   & $52$      & $T_{\text{ref}} > 0$ \\
$\kappa$ & Pénalité de défaut             & $0{,}15$  & $\kappa \gg \eta$ \\
$\delta$ & Dépréciation par inactivité    & $0{,}02$  & $\delta > 0$ \\
$\phi$   & Responsabilité Witness         & $0{,}15$  & $\phi \geq 3\eta$ \\
$\lambda$& Portabilité de base            & $0{,}20$  & $\lambda \in [0,1]$ \\
$\mu$    & Pénalité de migration          & $0{,}70$  & $\mu \in [0{,}6, 0{,}8]$ \\
$k_{\min}$ & Witnesses minimum            & $3$       & $k_{\min} \geq 3$ \\
$\alpha_k$ & Fraction adaptative          & $0{,}20$  & $\alpha_k \in (0,1)$ \\
\bottomrule
\end{tabular}
\end{table}

\subsection{Résultats}

Les résultats détaillés sont présentés à la Figure~\ref{fig:simulation}.

\begin{figure}[h]
\centering
\includegraphics[width=\textwidth]{ctp_analyse.png}
\caption{Résultats de simulation~: évolution des TrustScores sur 52 périodes,
scores finaux par comportement, analyse de sensibilité au paramètre $\delta$,
et vérification empirique des propositions théoriques. Les membres honnêtes
(vert) dominent strictement tous les autres comportements.}
\label{fig:simulation}
\end{figure}

\subsection{Vérification des propositions}

Les trois propositions théoriques sont vérifiées empiriquement~:

\begin{enumerate}
  \item \textbf{Bornitude~:} $\trustscore_{i,j}(t) \in [0,1]$ pour tous les
    membres, toutes les périodes. \hfill \checkmark

  \item \textbf{Dominance honnête sur le free-riding~:}
    $\bar{\trustscore}_{\text{honnête}} = 1{,}000 \gg
    \bar{\trustscore}_{\text{passager}} = 0{,}104$. \hfill \checkmark

  \item \textbf{Dominance honnête sur la collusion~:}
    $\bar{\trustscore}_{\text{honnête}} = 1{,}000 \gg
    \bar{\trustscore}_{\text{colludeur}} = 0{,}292$. \hfill \checkmark
\end{enumerate}

\subsection{Analyse de sensibilité}

L'analyse de sensibilité au paramètre $\delta$ confirme que pour
$\delta \in [0{,}01, 0{,}05]$, l'écart entre membres honnêtes et passagers
clandestins reste significatif. La valeur $\delta = 0{,}02$ offre un bon
équilibre entre punition de l'inactivité et tolérance aux absences
occasionnelles.

% ═════════════════════════════════════════════════════════
\section{Architecture des Nodes}
\label{sec:architecture}
% ═════════════════════════════════════════════════════════

\subsection{Hiérarchie à trois niveaux}

\CTP{} organise les communautés selon une hiérarchie de contraintes
descendantes~:

\begin{description}
  \item[Meta-Node $\mathcal{M}$] Définit les règles éthiques minimales et la
    structure mathématique du protocole. Immuable pour toutes les
    implémentations. Calcule $\rho_j$ et $\operatorname{sim}(j, j')$.

  \item[Node $N_j$] Définit ses paramètres dans les limites du Meta-Node.
    Peut être plus strict, jamais plus lâche. Gère ses propres Witnesses et
    son propre cycle de validation.

  \item[Sub-Node $N_{j,k}$] Sous-communauté thématique d'un Node. Hérite des
    contraintes du Node parent et peut en ajouter.
\end{description}

Un membre peut appartenir simultanément à plusieurs Nodes~; son TrustScore
est local à chaque Node et contribue à son score global via
l'équation~\eqref{eq:global}.

\subsection{Gouvernance interne}

Les décisions structurelles d'un Node (modification des paramètres, exclusion
d'un membre, changement du seuil $k_j$) sont traitées comme des
\textbf{GovSignals}~: des Signals spéciaux soumis au même protocole de
validation mais avec des exigences renforcées~:

\begin{itemize}
  \item Seuil de Witnesses $k_{\text{gov}} \geq 2 k_j$
  \item TrustScore minimum des Witnesses dans le quartile supérieur du Node
  \item Délai de contestation d'une période complète
\end{itemize}

Si un GovSignal est contesté par un nombre suffisant de membres, il est
suspendu et soumis à un vote pondéré par TrustScore.

\subsection{Mécanisme d'audit}

Tout membre peut émettre un \textbf{ChallengeSignal} contre une Contribution
existante. Ce mécanisme est incitatif~:

\begin{itemize}
  \item \textbf{Coût du challenge~:} fraction $\phi_c \cdot \trustscore_i$
    du score du challenger --- pour éviter les contestations abusives.
  \item \textbf{Si le challenge est accepté~:} la Contribution est invalidée,
    les Witnesses originaux sont pénalisés via $\phi$, le challenger récupère
    sa mise plus un bonus.
  \item \textbf{Si le challenge est rejeté~:} le challenger perd sa mise, les
    Witnesses ne sont pas affectés.
\end{itemize}

% ═════════════════════════════════════════════════════════
\section{Principes Éthiques Fondateurs}
\label{sec:ethique}
% ═════════════════════════════════════════════════════════

\CTP{} est un protocole technique porteur d'une philosophie précise. Les
principes suivants ne sont pas des recommandations~: ils sont constitutifs du
protocole. Toute implémentation qui les viole n'est pas conforme à \CTP{}.

\begin{description}
  \item[Autonomie communautaire] \CTP{} renforce la capacité d'une communauté
    à se gouverner sans dépendre d'une institution extérieure. Il ne remplace
    pas l'État~; il rend la communauté moins dépendante de lui pour ses
    fonctions de coordination interne.

  \item[Réciprocité] Chaque membre contribue et reçoit. Le parasitisme à long
    terme est mécaniquement impossible sous les paramètres recommandés.

  \item[Mémoire courte réparable] Un défaut passé ne condamne pas à vie.
    Le TrustScore peut se reconstruire~; cela prend du temps, proportionnel à
    la gravité du manquement.

  \item[Transparence des règles, confidentialité des personnes] Les règles du
    Node sont publiques et auditables. Les scores individuels ne le sont que
    si le Node en décide collectivement.

  \item[Non-marchandisation du score] Le TrustScore n'est ni un token ni un
    actif financier. Il ne fluctue pas selon un marché. Il ne s'échange pas.
    Toute tentative de monétisation directe du score contredit la philosophie
    fondatrice du protocole.

  \item[Décentralisation des données] Chaque Node gère ses propres données.
    \CTP{} ne centralise rien. Le Meta-Node calcule $\rho_j$ et
    $\operatorname{sim}(j,j')$ à partir de données agrégées et anonymisées.
\end{description}

% ═════════════════════════════════════════════════════════
\section{Complexité Algorithmique et Implémentation}
\label{sec:complexity}

\subsection{Analyse de complexité}

Le Tableau~\ref{tab:complexity} présente la complexité et les coûts gas
estimés sur Polygon pour chaque fonction principale du contrat
\texttt{CTPCore.sol}.

\begin{table}[h]
\centering
\caption{Complexité algorithmique et coûts gas estimés (Polygon mainnet).}
\label{tab:complexity}
\begin{tabular}{llll}
\toprule
\textbf{Fonction} & \textbf{Complexité} & \textbf{Gas estimé} & \textbf{Coût USD} \\
\midrule
\texttt{createNode()}      & $O(1)$    & $\sim 120{,}000$     & $< 0{,}02$ \\
\texttt{joinNode()}        & $O(1)$    & $\sim 80{,}000$      & $< 0{,}01$ \\
\texttt{submitSignal()}    & $O(1)$    & $\sim 70{,}000$      & $< 0{,}01$ \\
\texttt{validateSignal()}  & $O(n)$    & $\sim 50{,}000 + 5{,}000n$ & Variable \\
\texttt{advancePeriod()}   & $O(n)$    & $\sim 20{,}000n$     & Variable \\
\texttt{auditSignal()}     & $O(n)$    & $\sim 30{,}000n$     & Variable \\
\texttt{getTrustScore()}   & $O(1)$    & Gratuit (lecture)    & $0$ \\
\bottomrule
\end{tabular}
\end{table}

\subsection{Stratégie batch pour advancePeriod()}

La fonction \texttt{advancePeriod()} itère sur tous les membres du Node
--- complexité $O(n)$. Pour les Nodes de grande taille ($n > 50$),
le gas peut dépasser la limite par bloc. \CTP{} v1.1 recommande~:

\begin{itemize}
  \item Traitement par lots de 50 membres maximum~:
    \texttt{advancePeriod(nodeId, 0, 50)}, \texttt{advancePeriod(nodeId, 50, 50)}, etc.
  \item La période n'avance qu'une fois tous les membres traités.
  \item Coût total~: $\sim 20{,}000 \times n$ gas, soit $< 1{,}00$ USD
    pour $n = 200$ membres sur Polygon.
\end{itemize}

\section{Travaux Connexes et Positionnement}
\label{sec:related}
% ═════════════════════════════════════════════════════════

\subsection{Confiance computationnelle --- Marsh (1994)}

La première formalisation mathématique de la confiance computationnelle est
due à Marsh~\cite{marsh1994trust}, dont la thèse de 1994 introduit le
concept de confiance comme variable continue entre 0 et 1, avec des
propriétés de mise à jour basées sur l'expérience. \CTP{} étend ce cadre
fondateur en introduisant~: (1) la validation communautaire collective via
les Witnesses, (2) la hiérarchie de Nodes, et (3) l'ancrage dans la
philosophie Ubuntu et la logique des tontines africaines.

\subsection{Systèmes de réputation pair-à-pair}

Les travaux fondateurs sur EigenTrust~\cite{kamvar2003eigentrust}
et PeerTrust~\cite{xiong2004peertrust} mesurent la fiabilité
des nœuds dans des réseaux P2P à partir de l'historique des transactions.
\CTP{} s'en distingue sur trois points~: (1) l'introduction d'une
hiérarchie de Nodes paramétrable, (2) le mécanisme de responsabilité des
Witnesses qui engage la réputation du validateur, et (3) l'ancrage dans une
philosophie de confiance communautaire africaine plutôt que dans un modèle
purement transactionnel.

\subsection{Gouvernance décentralisée et DAOs}

Les organisations autonomes décentralisées (DAOs) confondent généralement
richesse et gouvernance~: le pouvoir de vote est proportionnel au nombre de
tokens détenus. \CTP{} propose une alternative radicale où le pouvoir de vote
est proportionnel au TrustScore --- une mesure comportementale non monétisable.
Cette approche rejoint les préoccupations soulevées par Buterin sur les
limites des systèmes de gouvernance token-pondérés~\cite{buterin2021governance}.

\subsection{Modélisation formelle de la négociation dans les espaces de données}

Le travail récent de Daoudi, Lefrançois et Zimmermann~\cite{daoudi2026negotiation}
propose un cadre formel pour la négociation automatisée des contrats d'accès
aux données dans les espaces de données décentralisés. Leur approche identifie
explicitement l'absence d'un mécanisme garantissant la fiabilité des acteurs
au-delà de la négociation elle-même comme un verrou scientifique majeur.

\CTP{} répond précisément à ce manque~: le TrustScore fournit une mesure
continue de la fiabilité comportementale, le mécanisme de Witnesses assure
un audit communautaire continu des engagements, et la théorie des jeux
garantit la résistance aux comportements stratégiques post-contractuels.

Inversement, \CTP{} spécifie une fonction de similarité sémantique
$\operatorname{sim}(j, j')$ entre Nodes --- nécessaire pour la portabilité
inter-communautés --- sans la formaliser complètement. Le cadre ODRL et les
ontologies de domaine développés par Daoudi et al. constituent des outils
naturels pour cette formalisation. La combinaison des deux approches ouvre
la perspective d'un protocole d'échange décentralisé complet~: le cadre
de Daoudi et al. comme moteur de négociation, \CTP{} comme système de
garantie de la réputation des acteurs.

\subsection{Microfinance et tontines}

Les travaux empiriques sur les tontines africaines~\cite{ardener1964rotating}
et les groupes d'entraide~\cite{besley1993rotating} documentent l'efficacité
des mécanismes de pression sociale et de réputation collective comme
substituts aux garanties bancaires formelles. \CTP{} formalise ces mécanismes
dans un cadre mathématique rigoureux et les rend implémentables dans des
systèmes numériques décentralisés.

% ═════════════════════════════════════════════════════════
\section{Perspectives}
\label{sec:conclusion}
% ═════════════════════════════════════════════════════════

\subsection{Limites actuelles}

Ce document est une spécification fondatrice --- il pose des bases
théoriques solides mais comporte des limites explicites~:

\begin{description}
  \item[sim(j,j') -- portabilité partielle] La formalisation directionnelle
    de v1.1 est calculable on-chain mais ne capture pas la proximité
    sémantique entre domaines sans historique commun. L'intégration des
    ontologies ODRL de Daoudi et al.~\cite{daoudi2026negotiation} est
    envisagée en v1.2.

  \item[Calibration empirique] Les paramètres $\eta$, $\kappa$, $\delta$
    sont validés par simulation mais non calibrés sur des données réelles
    de tontines. Une collecte de données en Afrique de l'Ouest est
    nécessaire pour la v2.0.

  \item[Attaque Sybil d'élite] Le protocole résiste aux faux comptes de
    masse via le mécanisme de Witnesses. Un groupe de membres établis en
    collusion reste une menace partiellement couverte par $\phi$ et le
    ChallengeSignal.

  \item[Scalabilité on-chain] \texttt{advancePeriod()} est $O(n)$.
    Au-delà de 200 membres par Node, la stratégie batch est obligatoire.
\end{description}

\subsection{Ce que ce document établit}

Ce document constitue la spécification fondatrice de \CTP{}~:
une formalisation mathématique rigoureuse des mécanismes de confiance
communautaire, une analyse par la théorie des jeux garantissant la résistance
stratégique, et une validation empirique par simulation.

\subsection{Travaux futurs}

\paragraph{Implémentation en smart contracts (Phase 4).}
La traduction du moteur \CTP{} en Solidity est la prochaine étape. La
structure mathématique a été conçue dès l'origine pour être implémentable
on-chain~: toutes les opérations sont déterministes et calculables en temps
polynomial.

\paragraph{Calibration sur données réelles (Phase 5).}
Les paramètres recommandés sont fondés sur la structure théorique et validés
par simulation. Une calibration sur des données réelles de tontines (Afrique
de l'Ouest) et de coopératives permettrait d'affiner $\eta$, $\kappa$ et
$\delta$.

\paragraph{Distance sémantique inter-Nodes.}
La fonction $\operatorname{sim}(j, j')$ est spécifiée mais non encore
implémentée. Une approche par plongements sémantiques (\textit{embeddings}) des
descriptions de Nodes est envisagée.

\paragraph{Extension aux réseaux sociaux.}
L'application de \CTP{} aux réseaux sociaux communautaires requiert la
définition formelle de ce qu'est un Signal de qualité dans ce contexte~:
une question ouverte qui nécessite une collaboration avec des communautés
existantes.

\subsection{Invitation à contribuer}

\CTP{} est un protocole ouvert. OpenScience Community invite toute
communauté, chercheur ou développeur à~:

\begin{itemize}
  \item Tester le protocole sur des cas d'usage concrets
  \item Proposer des modifications formalisées via le mécanisme de GovSignal
  \item Contribuer à l'implémentation open source
  \item Partager des données de calibration
\end{itemize}

\medskip
\begin{center}
\textit{``Je suis parce que nous sommes.''}\\
\medskip
{\small CommunityTrust Protocol v0.1 --- OpenScience Community\\
Document vivant --- Creative Commons BY-SA 4.0}
\end{center}

% ═════════════════════════════════════════════════════════
\appendix
\section{Algorithme de mise à jour du TrustScore}
% ═════════════════════════════════════════════════════════

\begin{algorithm}[H]
\caption{Mise à jour du TrustScore --- CTP v0.1}
\begin{algorithmic}[1]
\Require Membres $\mathcal{M}_j$, Signals validés $\mathcal{S}(t)$, Défauts $\mathcal{D}(t)$, Actifs $\mathcal{A}(t)$
\Ensure TrustScores mis à jour $\trustscore_{i,j}(t+1)$

\ForAll{$s \in \mathcal{S}(t)$}
  \If{$q(s) < q_{\min}$ \textbf{et} $\text{Uniforme}(0,1) < p_{\text{audit}}$}
    \State $s.\text{audité} \leftarrow \top$
  \EndIf
\EndFor

\ForAll{$i \in \mathcal{M}_j$}
  \State $G \leftarrow \eta \cdot \sum_{s \in \mathcal{S}(t) : s.\text{auteur}=i,\; \neg s.\text{audité}} q(s)$
  \State $L \leftarrow \kappa \cdot \1_{[i \in \mathcal{D}(t)]}$
  \State $D \leftarrow \delta \cdot \trustscore_{i,j}(t) \cdot \1_{[i \notin \mathcal{A}(t)]}$
  \State $W \leftarrow \phi \cdot \sum_{s : i \in \mathcal{W}(s),\; s.\text{audité}} \frac{\trustscore_{i,j}(t)}{|\mathcal{W}(s)|}$
  \State $\trustscore_{i,j}(t+1) \leftarrow \clip(\trustscore_{i,j}(t) + G - L - D - W,\; 0,\; 1)$
\EndFor
\end{algorithmic}
\end{algorithm}

\section{Table des symboles}

\begin{table}[H]
\centering
\caption{Table complète des symboles utilisés dans ce document.}
\begin{tabular}{cll}
\toprule
\textbf{Symbole} & \textbf{Signification} & \textbf{Domaine} \\
\midrule
$\trustscore_{i,j}(t)$ & TrustScore du membre $i$ dans Node $j$ au temps $t$ & $[0,1]$ \\
$\trustscore_i^{\mathcal{M}}(t)$ & TrustScore Global du membre $i$ & $[0,1]$ \\
$G_{i,j}(t)$ & Gain par contribution & $\R_+$ \\
$L_{i,j}(t)$ & Perte par défaut & $\R_+$ \\
$D_{i,j}(t)$ & Dépréciation par inactivité & $\R_+$ \\
$W_{i,j}(t)$ & Coût de mauvaise validation & $\R_+$ \\
$q(s)$ & Qualité d'un Signal validé & $[0,1]$ \\
$\mathcal{W}(s)$ & Ensemble des Witnesses du Signal $s$ & --- \\
$k_j$ & Seuil adaptatif de Witnesses & $\N$ \\
$\omega_{i,j}(t)$ & Poids du Node $j$ pour le membre $i$ & $[0,1]$ \\
$\rho_j$ & Indice de difficulté du Node $j$ & $\R_+$ \\
$\lambda_{\text{eff}}$ & Portabilité effective & $[0,1]$ \\
$\operatorname{sim}(j,j')$ & Similarité sémantique entre Nodes & $[0,1]$ \\
$V_{i,j}(t)$ & Pouvoir de vote du membre $i$ & $[0,1]$ \\
$\eta$ & Taux de gain & $(0,1)$ \\
$\kappa$ & Pénalité de défaut & $(0,1)$ \\
$\delta$ & Taux de dépréciation & $(0,1)$ \\
$\phi$ & Coefficient de responsabilité & $(0,1)$ \\
$\lambda$ & Portabilité de base & $(0,1)$ \\
$\mu$ & Pénalité de migration & $(0,1)$ \\
$\beta$ & Facteur d'escompte temporel & $(0,1)$ \\
\bottomrule
\end{tabular}
\end{table}

% ═════════════════════════════════════════════════════════
\begin{thebibliography}{99}
% ═════════════════════════════════════════════════════════

\bibitem{marsh1994trust}
S.~P. Marsh,
\textit{Formalising Trust as a Computational Concept},
Thèse de doctorat, University of Stirling, 1994.
Première formalisation mathématique de la confiance computationnelle~;
ancêtre intellectuel direct de \CTP{}.

\bibitem{kamvar2003eigentrust}
S.~D. Kamvar, M.~T. Schlosser, et H.~Garcia-Molina,
\textit{The EigenTrust algorithm for reputation management in P2P networks},
Proceedings of the 12th International Conference on World Wide Web (WWW),
pp.~640--651, 2003.

\bibitem{xiong2004peertrust}
L.~Xiong et L.~Liu,
\textit{PeerTrust: Supporting Reputation-Based Trust for Peer-to-Peer Electronic Communities},
IEEE Transactions on Knowledge and Data Engineering,
vol.~16, no.~7, pp.~843--857, 2004.

\bibitem{buterin2021governance}
V.~Buterin,
\textit{Moving beyond coin voting governance},
vitalik.ca, août 2021.
Disponible~: \url{https://vitalik.ca/general/2021/08/16/voting3.html}

\bibitem{daoudi2026negotiation}
K.~Daoudi, M.~Lefrançois, et A.~Zimmermann,
\textit{Vers une modélisation formelle de la négociation des contrats d'accès
aux données dans les espaces de données},
Mines Saint-Étienne, Univ. Clermont Auvergne, CNRS, UMR 6158 LIMOS,
mars 2026.

\bibitem{ardener1964rotating}
S.~Ardener,
\textit{The comparative study of rotating credit associations},
Journal of the Royal Anthropological Institute of Great Britain and Ireland,
vol.~94, no.~2, pp.~201--229, 1964.

\bibitem{besley1993rotating}
T.~Besley, S.~Coate, et G.~Loury,
\textit{The Economics of Rotating Savings and Credit Associations},
The American Economic Review,
vol.~83, no.~4, pp.~792--810, 1993.

\bibitem{ostrom1990governing}
E.~Ostrom,
\textit{Governing the Commons: The Evolution of Institutions for Collective Action},
Cambridge University Press, 1990.

\bibitem{nakamoto2008bitcoin}
S.~Nakamoto,
\textit{Bitcoin: A Peer-to-Peer Electronic Cash System},
2008.
Disponible~: \url{https://bitcoin.org/bitcoin.pdf}

\end{thebibliography}

\end{document}
